\documentclass[11pt]{article}

\usepackage[a4paper,margin=1in]{geometry}
\usepackage{amsmath,amssymb,bm}
\usepackage{booktabs}
\usepackage{hyperref}

\hypersetup{
    colorlinks=true,
    linkcolor=blue,
    citecolor=blue,
    urlcolor=blue
}

\newcommand{\eps}{\varepsilon}
\newcommand{\mean}[1]{\langle #1 \rangle}
\newcommand{\master}{geqoe\_averaged\_j2\_zonal\_note}

\title{GEqOE-Enhanced Conjunction Filter\\
\large Replacing the First-Order Radius Model in the SO-Filter\\
with Exact-in-Eccentricity Analytical Reconstruction}
\author{Jose Manuel Montilla Garc\'ia}
\date{\today}

\begin{document}
\maketitle

\begin{abstract}
This companion note describes how the GEqOE averaged zonal theory (master
note, \texttt{\master{}}) can be integrated into the Space Occupancy (SO)
conjunction filter framework of Rivero, Bombardelli, and
Vazquez~\cite{rivero2025}. The SO-filter represents a major advance over the
classical apogee-perigee (AP) filter for first-stage conjunction screening,
reducing false positives by two orders of magnitude while maintaining linear
computational scaling. However, its current formulation is restricted to
orbits with eccentricity $e < 0.1$ due to first-order eccentricity expansions
in the radius model and Cook's circle theory. The GEqOE averaged theory
removes this restriction by providing: (i)~an exact-in-eccentricity osculating
radius reconstruction via the short-period map, (ii)~a replacement of the
Kozai--Lyddane osculating-to-mean conversion with the GEqOE first-order
inverse map, (iii)~inclusion of zonal harmonics $J_2$--$J_5$ instead of
$J_2$-only, and (iv)~a consistent drag extension. The estimated improvement
is a factor 2--5 reduction in the required correction buffers and extension of
the SO-filter to the full catalogued population including high-eccentricity
orbits such as GTO and Molniya.
\end{abstract}

\tableofcontents

%% ====================================================================
\section{Motivation and Context}
\label{sec:motivation}

\subsection{The conjunction screening problem}

Conjunction analysis (CA) for the catalogued resident space object (RSO)
population involves screening $\sim\!10^{11}$ pairs over typical 5-day
windows. First-stage filtering reduces this set to a manageable number of
candidates ($\sim\!10^4$--$10^5$) for detailed orbit determination and
probability-of-collision computation. The classical apogee-perigee (AP) filter
eliminates pairs whose radial extent bands do not overlap. As shown by
Rivero et al.~\cite{rivero2025}, the AP-filter produces $\sim\!224{,}000$
false positives from $\sim\!144$ million pairs (16,972 LEO orbits), requiring
buffers of 3--6~km over 5-day windows to avoid false negatives.

\subsection{The SO-filter advance}

The SO-filter~\cite{rivero2025} replaces the simple apogee/perigee bounds
with a refined radial occupancy model based on Cook's frozen orbit
theory~\cite{cook1966} and short-term extensions via a quartic equation. This
reduces false positives from $224{,}316$ (AP) to $2{,}242$ (SO)---a
\textbf{100-fold improvement}---with only a 15\% increase in computation time
(8.30~min vs.\ 7.20~min for 16,972 orbits). The required correction buffers
drop from 3--6~km (AP) to 0.5--1.2~km (SO) for zero false negatives.

\subsection{Current limitations}

The SO-filter's accuracy rests on three analytical approximations, each of
which the GEqOE theory can improve:
\begin{enumerate}
\item \textbf{First-order eccentricity expansion}: The radius polar equation
(Eq.~19 in~\cite{rivero2025}) expands $r(\hat\theta, \beta)$ to first order
in eccentricity, limiting applicability to $e < 0.1$. High-eccentricity
orbits (GTO, Molniya, HEO) are excluded entirely.
\item \textbf{$J_2$-only physics}: Cook's circle theory and the frozen
eccentricity $e_f$ (Eq.~4 in~\cite{rivero2025}) include only $J_2$ and odd
zonals at first order. The actual geopotential through $J_5$ causes additional
frozen-eccentricity shifts and secular rate corrections.
\item \textbf{Kozai--Lyddane mean elements}: The osculating-to-mean
conversion (Appendix~B of~\cite{rivero2025}) uses Kozai--Lyddane theory, which
is singular at the critical inclination ($i \approx 63.4^\circ$) and
introduces small but nonzero errors that contribute to the correction buffer.
\end{enumerate}

The GEqOE averaged theory addresses all three limitations simultaneously.

%% ====================================================================
\section{Component-by-Component Replacement}
\label{sec:components}

This section maps each analytical component of the SO-filter to its GEqOE
replacement, identifying exactly where each piece enters the computation
pipeline.

\subsection{Osculating-to-mean conversion (replacing Kozai--Lyddane)}
\label{sec:osc2mean}

\subsubsection{Current approach}

The SO-filter converts initial osculating orbital elements
$\{a, e, i, \omega, \Omega, M\}$ to mean elements
$\{\hat a, \hat e, \hat i, \hat\omega, \hat\Omega, \hat M\}$ using the
Kozai~\cite{kozai1959}--Lyddane~\cite{lyddane1963} first-order $J_2$
short-period removal (Appendix~B of~\cite{rivero2025}, Eqs.~69--87). This
conversion:
\begin{itemize}
\item accounts for $J_2$ only;
\item is first-order in $J_2$ (errors of order $J_2^2 \sim 10^{-6}$);
\item contains a $1/\sin i$ singularity near equatorial orbits and a
$1/(1-5\cos^2 i)$ singularity near the critical inclination;
\item requires solving Kepler's equation (Eq.~69 of~\cite{rivero2025}).
\end{itemize}

\subsubsection{GEqOE replacement}

The GEqOE first-order inverse map (master note, Section~8) converts
osculating Cartesian state vectors $(\bm r, \bm v)$ to mean GEqOE
$\bar{\bm y} = (\bar\nu, \bar p_1, \bar p_2, \bar K, \bar q_1, \bar q_2)$
via:
\begin{equation}
\bar{\bm y} = \bm y_{\rm osc} - \eps\,\bm u_1(\bm y_{\rm osc}),
\label{eq:inverse_map}
\end{equation}
where $\bm u_1$ is the same short-period correction evaluated at the
osculating state. This map:
\begin{itemize}
\item includes $J_2$--$J_5$ (all zonals in the theory);
\item is exact in eccentricity---no Taylor expansion in $e$;
\item is free of the equatorial and critical-inclination singularities of the
classical elements (though singular for retrograde equatorial and rectilinear
orbits, which are not relevant for the RSO catalogue);
\item requires no Kepler equation solve---the GEqOE state uses $K$ directly.
\end{itemize}

\subsubsection{Impact on the SO-filter pipeline}

The osculating-to-mean conversion is the \emph{entry point} of the SO-filter:
every RSO in the catalogue passes through it once at the reference epoch $t_0$.
The GEqOE inverse map is algebraic (rational functions of $q$ and $Q$,
trigonometric functions of $K$), with comparable computational cost to
Kozai--Lyddane. The improvement is:
\begin{itemize}
\item elimination of the $\sin i$ and critical-inclination singularities,
\item inclusion of $J_3$--$J_5$ short-period effects (up to $\sim\!100$~m for
low-altitude orbits),
\item exact-in-$e$ short-period removal, reducing the initialization error
that propagates into the radial bounds.
\end{itemize}

The validated round-trip error of the GEqOE inverse map is 7.2~m
(low-$e$) and 1.6~m (high-$e$) at first order (master note, Section~9),
compared to typical Kozai--Lyddane errors of $\sim\!100$~m for $J_2$-only.

\subsection{Mean eccentricity evolution (replacing Cook's circle)}
\label{sec:cook_replacement}

\subsubsection{Current approach}

Cook's circle theory~\cite{cook1966} describes the evolution of the mean
eccentricity vector $(\xi, \eta)$ as a circle of radius $e_p$ centered at
$(0, e_f)$ in the nodal frame (Eqs.~1--9 of~\cite{rivero2025}):
\begin{equation}
\xi(\tau) = e_p \cos\beta(\tau), \qquad
\eta(\tau) = e_p \sin\beta(\tau) + e_f,
\label{eq:cook_circle}
\end{equation}
where $e_f$ is the frozen eccentricity (Eq.~4), $e_p$ the proper
eccentricity (Eq.~5), and $\beta(\tau) = k\tau + \alpha$ with precession rate
$k$ (Eq.~9). The theory is:
\begin{itemize}
\item first-order in eccentricity;
\item $J_2$-only for the precession rate $k$, with higher odd zonals appearing
only in $e_f$;
\item valid for ``small'' eccentricity (the paper applies it to $e < 0.1$).
\end{itemize}

\subsubsection{GEqOE replacement}

The GEqOE mean ODE system (master note, Section~5) provides the exact
first-order mean rates for all six mean elements $\bar{\bm y}$. For the
pure $J_2$ case, the mean system is especially simple (master note,
Section~5.1):
\begin{itemize}
\item $\bar\nu$, $\bar g$, $\bar Q$ are \emph{exactly constant};
\item $(\bar p_1, \bar p_2)$ rotate uniformly at rate $\dot\omega_{J_2}$;
\item $(\bar q_1, \bar q_2)$ rotate uniformly at rate $\dot\Omega_{J_2}$.
\end{itemize}
For the general zonal case ($J_2$--$J_5$), the mean ODE is a smooth,
slowly varying system that can be integrated cheaply (or via Taylor polynomial
maps as described in the companion note \texttt{geqoe\_mean\_taylor\_prop}).

The eccentricity magnitude $g = |\bm g|$ and the argument of pericenter
embedded in $(p_1, p_2)$ are \emph{implicit} in the GEqOE formulation; the
Cook-circle behavior is recovered as a \emph{consequence} of the mean ODE
rather than an \emph{assumption}. Crucially, the GEqOE mean drift:
\begin{itemize}
\item is exact in eccentricity (no first-order-in-$e$ truncation);
\item includes cross-zonal effects from $J_3$, $J_4$, $J_5$ on both the
precession rate and the frozen point;
\item does not require the ``small eccentricity'' assumption.
\end{itemize}

\subsubsection{Impact on the SO-filter}

The mean eccentricity evolution feeds directly into the radius polar equation
(Section~\ref{sec:radius_replacement}). More accurate mean elements at each
evaluation point $\tau$ reduce the error in $r(\hat\theta, \beta)$ and thus
the required correction buffer. The improvement is most significant for:
\begin{itemize}
\item orbits near the critical inclination, where $J_3$--$J_5$ contributions
to the frozen eccentricity become comparable to the $J_2$ term;
\item moderate-eccentricity orbits ($0.01 < e < 0.1$), where first-order-in-$e$
errors in Cook's circle accumulate over the 5-day screening window;
\item high-eccentricity orbits ($e > 0.1$), which are currently excluded
entirely.
\end{itemize}

\subsection{Radius reconstruction (replacing the first-order polar equation)}
\label{sec:radius_replacement}

\subsubsection{Current approach}

The SO-filter's radius polar equation (Eq.~19 of~\cite{rivero2025}) is:
\begin{equation}
r \approx \hat a\bigl(1 - \hat e\cos(\hat\theta - \hat\omega)\bigr)
+ \frac{J_2}{4\hat a}\bigl[(9 + \cos 2\hat\theta)\sin^2\hat\imath - 6\bigr],
\label{eq:rivero_radius}
\end{equation}
which is first-order in eccentricity (the $1 - e\cos f$ term) and
first-order in $J_2$ (the $J_2/(4\hat a)$ correction). This formula is the
heart of both the long-term SO bounds (Eqs.~13--14) and the short-term SO
quartic equation (Eq.~34).

The formula's error has two sources:
\begin{enumerate}
\item \textbf{Eccentricity truncation}: The exact Keplerian radius is
$r = a(1-e^2)/(1+e\cos f)$, which differs from $a(1-e\cos f)$ by
$O(e^2)$. For $e = 0.1$ this is $\sim\!1\%$ of $a$ ($\sim\!70$~km for LEO),
though the \emph{mean} error over the orbit is smaller.
\item \textbf{$J_2$-only geopotential}: The $J_2/(4\hat a)[\cdots]$ term
misses $J_3$--$J_5$ contributions which can shift the extremal radii by
tens of meters.
\end{enumerate}

\subsubsection{GEqOE replacement: exact osculating reconstruction}

The GEqOE short-period map (master note, Sections~6--7) provides:
\begin{equation}
\bm y_{\rm osc} = \bar{\bm y} + \eps\,\bm u_1(\bar{\bm y}, \bar K),
\label{eq:forward_map}
\end{equation}
where $\bm u_1$ is known in closed form (polynomial-plus-residue expressions
in $q$ and $Q$, rational in $F = e^{if}$). From the osculating GEqOE state,
the osculating radius is recovered \emph{exactly} via the GEqOE-to-Cartesian
transformation (master note, Section~3):
\begin{equation}
r = \frac{c^2}{\mu}\frac{1}{1 + g\cos f},
\label{eq:exact_radius}
\end{equation}
where $c$, $g$, and $f$ are extracted from the osculating GEqOE state.
This reconstruction:
\begin{itemize}
\item is \emph{exact in eccentricity}---no Taylor expansion in $e$;
\item includes $J_2$--$J_5$ in the short-period correction;
\item captures the correct osculating radius at every point on the orbit, not
just an approximation valid for small $e$.
\end{itemize}

\subsubsection{Two implementation strategies}

There are two ways to use the GEqOE reconstruction within the SO-filter
framework:

\paragraph{Strategy A: Direct replacement of the radius polar equation.}
For a given mean state $\bar{\bm y}$ and a grid of $\bar K$ values spanning
$[0, 2\pi]$, evaluate the forward map~\eqref{eq:forward_map} and compute
$r(\bar K)$ via~\eqref{eq:exact_radius}. The maximum and minimum $r$ over
the grid give the short-term SO bounds. This replaces the quartic equation
(Eq.~34 of~\cite{rivero2025}) with a direct sweep.

The computational cost is $N_{\rm grid}$ evaluations of the short-period map
per RSO per evaluation epoch. With $N_{\rm grid} = 36$ (every $10^\circ$) and
algebraic evaluation of the short-period map, this is comparable to solving
the quartic equation and evaluating the critical-point Hessian.

\paragraph{Strategy B: Improved quartic equation.}
Retain the quartic structure of the SO-filter but replace the first-order
radius formula~\eqref{eq:rivero_radius} with a higher-fidelity expression
derived from the GEqOE reconstruction. Specifically, expand
$r(\bar{\bm y}, \bar K)$ as a truncated Fourier series in $\hat\theta$ and
$\beta$:
\begin{equation}
r(\hat\theta, \beta) \approx r_0 + \sum_{k=1}^{N_h}
\bigl[A_k \cos k\hat\theta + B_k \sin k\hat\theta\bigr]
+ \text{zonal corrections},
\label{eq:improved_radius}
\end{equation}
where the coefficients $r_0$, $A_k$, $B_k$ are computed from the GEqOE
short-period map at the current mean state. Setting $dr/d\hat\theta = 0$ and
$dr/d\beta = 0$ yields a polynomial system whose degree depends on $N_h$.
For $N_h = 2$ the polynomial is quartic (same as the current SO-filter); for
$N_h = 3$ it is sextic, solvable via companion-matrix eigenvalues.

Strategy~A is simpler to implement and more robust (no root-finding
pathologies); Strategy~B preserves the quartic structure for backward
compatibility but may gain less in accuracy. We recommend Strategy~A for new
implementations and Strategy~B for incremental upgrades of existing code.

\subsubsection{Estimated accuracy improvement}

The SO-filter's current model error is $\sim\!0.5$~km mean, with 98.7\% of
cases below 1~km (Fig.~6 of~\cite{rivero2025}). The GEqOE reconstruction
error is:
\begin{itemize}
\item $\sim\!0.1$~km ($\sim\!100$~m) for low-$e$ orbits (master note,
Section~9: sub-kilometer osculating reconstruction);
\item $\sim\!1$--$3$~km for $e = 0.65$ (first-order theory), but this regime
is entirely \emph{excluded} from the current SO-filter.
\end{itemize}
For the $e < 0.1$ population where both models apply, the GEqOE
reconstruction should reduce the model error by a factor of $\sim\!3$--$5$
(from $\sim\!0.5$~km to $\sim\!0.1$~km), because:
\begin{enumerate}
\item the eccentricity truncation error ($O(e^2 a) \sim 70$~m for $e = 0.1$)
is eliminated;
\item the $J_3$--$J_5$ short-period corrections ($\sim\!50$~m) are included;
\item the improved inverse map reduces initialization error from
$\sim\!100$~m to $\sim\!7$~m.
\end{enumerate}

\subsection{Buffer reduction}
\label{sec:buffer}

The correction buffer is the safety margin added to each RSO's radial bounds
to eliminate false negatives. Rivero et al.\ report (Table~3
of~\cite{rivero2025}):
\begin{center}
\begin{tabular}{lcc}
\toprule
& AP-filter [km] & SO-filter [km] \\
\midrule
Low-$e$, low altitude ($< 400$~km) & 11.5 & 0.98 \\
Low-$e$, mid altitude (400--700~km) & 11.3 & 1.28 \\
Low-$e$, high altitude ($> 1000$~km) & 8.6 & 2.03 \\
Moderate-$e$, low altitude & 10.7 & 0.90 \\
Moderate-$e$, high altitude & 8.5 & 2.51 \\
\bottomrule
\end{tabular}
\end{center}

The buffer is driven by the worst-case model error over the apsidal
precession period (Eq.~43 of~\cite{rivero2025}):
\begin{equation}
\varepsilon^* = \max\{r^{\rm short}_{\max} - r^{\rm num}_{\max},\;
r^{\rm num}_{\min} - r^{\rm short}_{\min}\}.
\end{equation}
Since the GEqOE reconstruction reduces the model error by a factor of
$\sim\!3$--$5$ for the $e < 0.1$ population, the required buffers should
shrink proportionally:
\begin{center}
\begin{tabular}{lcc}
\toprule
& SO-filter [km] & GEqOE-SO [km] (est.) \\
\midrule
Low-$e$, low altitude & 0.98 & 0.2--0.3 \\
Low-$e$, mid altitude & 1.28 & 0.3--0.4 \\
Low-$e$, high altitude & 2.03 & 0.4--0.7 \\
Moderate-$e$ & 0.90--2.51 & 0.2--0.5 \\
High-$e$ ($e > 0.1$) & N/A & 1--3 \\
\bottomrule
\end{tabular}
\end{center}

The last row is new: the GEqOE-enhanced SO-filter would for the first time
provide analytical radial bounds for high-eccentricity orbits, which
currently require numerical propagation or are excluded from the filter
entirely. Even with larger buffers (1--3~km for $e > 0.1$), this is
substantially better than the AP-filter's 8--12~km buffers.

\subsubsection{Impact on false positive rates}

Smaller buffers translate directly into fewer false positives. The SO-filter
with category-specific buffers achieves $\rho_{\rm FP} = 1.68\%$ (Table~4
of~\cite{rivero2025}). Since the number of false positives scales roughly as
the square of the buffer size (two-dimensional radial band overlap), a factor
$\sim\!3$ reduction in buffer size implies a factor $\sim\!9$ reduction in
false positives:
\begin{equation}
\rho_{\rm FP}^{\rm GEqOE} \sim \frac{\rho_{\rm FP}^{\rm SO}}{9}
\sim 0.2\%.
\end{equation}
This would bring the false positive count from $\sim\!537{,}000$ (SO-filter
with buffer) to $\sim\!60{,}000$---a further $\sim\!9\times$ reduction on top
of the SO-filter's already dramatic improvement over the AP-filter.

%% ====================================================================
\section{Eccentricity Range Extension}
\label{sec:eccentricity}

\subsection{The current restriction}

Rivero et al.\ explicitly state: ``it is important to underline that orbits
with eccentricities higher than 0.1 and apogee radii exceeding 40,000~km
have been left out.'' This excludes:
\begin{itemize}
\item Geostationary transfer orbits (GTO): $e \approx 0.73$
\item Molniya orbits: $e \approx 0.74$
\item Highly eccentric orbits (HEO): $e \approx 0.4$--$0.8$
\item GPS transfer orbits: $e \approx 0.02$--$0.1$
\end{itemize}

These orbit classes contain significant debris populations and active
satellites. Including them in the first-stage filter (rather than requiring
numerical propagation) would reduce the computational burden on the
downstream conjunction assessment pipeline.

\subsection{Why GEqOE removes the restriction}

The GEqOE averaged theory uses the eccentricity parameter
$q = g/(1 + \beta)$ where $\beta = \sqrt{1 - g^2}$ and $g$ is the
GEqOE eccentricity magnitude. The entire theory---mean ODE, short-period
map, and inverse map---is formulated as rational functions of $q$ and the
complex true anomaly exponential $F = e^{if}$, with \emph{no Taylor
expansion in eccentricity}. The only restriction is $g < 1$ (non-rectilinear
orbits) and $q_1^2 + q_2^2 < 1$ (non-retrograde-equatorial).

Validation in the master note (Section~9) demonstrates:
\begin{itemize}
\item Position reconstruction error $< 1$~km for $e = 0.01$ (10 orbits);
\item Position reconstruction error $\sim\!3$~km for $e = 0.65$ (10 orbits);
\item Scaling slope $\approx 1.0$ (first-order theory confirmed).
\end{itemize}

For the SO-filter application, the relevant quantity is the \emph{radial}
component of the reconstruction error, which is typically smaller than the
full 3D position error (the dominant error is usually in the along-track
direction). We estimate:
\begin{equation}
\delta r_{\rm GEqOE} \sim
\begin{cases}
50\text{--}100~\text{m} & e < 0.1, \\
200\text{--}500~\text{m} & 0.1 < e < 0.5, \\
0.5\text{--}2~\text{km} & e > 0.5.
\end{cases}
\end{equation}

\subsection{Population coverage}

Including orbits up to $e = 0.8$ would extend the SO-filter from the current
16,972 orbits (restricted catalogue) to the full catalogued LEO/MEO/GTO
population of $\sim\!25{,}000$+ objects. The high-eccentricity orbits are
precisely those for which the AP-filter performs worst (large radial range),
so even moderate GEqOE-SO accuracy provides a significant advantage.

%% ====================================================================
\section{Drag Integration}
\label{sec:drag}

\subsection{Current drag treatment in the SO-filter}

Rivero et al.\ incorporate atmospheric drag via the Billik--Karrenberg
altitude decay formula (Eq.~45 of~\cite{rivero2025}):
\begin{equation}
h(t) = \frac{1}{\beta_a}\log\Bigl[e^{\beta_a h_0}
- B\sqrt{\mu R_\oplus \beta_a \bar\rho}\; t\Bigr],
\label{eq:billik}
\end{equation}
where $B = C_D A_c / m_c$ is the ballistic coefficient, $\beta_a$ and
$\bar\rho$ are exponential density model parameters (Table~5
of~\cite{rivero2025}), and $h_0$ is the initial altitude. This formula:
\begin{itemize}
\item is valid for near-circular orbits only;
\item treats drag as a secular altitude decay (no eccentricity evolution);
\item adds a correction of $\sim\!0.6$~km to the minimum radius bounds for
altitudes below 500~km.
\end{itemize}

\subsection{GEqOE drag extension}

The companion note \texttt{geqoe\_mean\_drag\_extension} develops a more
complete drag model within the GEqOE framework. The key features are:
\begin{itemize}
\item Secular drag rates for both $\bar\nu$ (energy decay) and $\bar g$
(eccentricity circularization) via King-Hele formulas;
\item Bessel-function expansion of the exponential density model, valid for
arbitrary eccentricity;
\item Addition to the mean ODE without rederiving the zonal short-period map
(justified because drag short-period effects are $O(\eps_{\rm drag}
\cdot \eps_{\rm zonal})$).
\end{itemize}

This provides a natural upgrade path: the GEqOE-enhanced SO-filter inherits
the drag extension automatically through the mean ODE, rather than requiring
a separate altitude-decay formula. The improvement is especially relevant for:
\begin{itemize}
\item eccentric orbits below 500~km perigee, where drag circularization is
significant over 5-day windows;
\item very low orbits ($< 300$~km), where the Billik formula's circular-orbit
assumption introduces errors.
\end{itemize}

%% ====================================================================
\section{Computational Architecture}
\label{sec:computational}

\subsection{Scaling requirements}

The ``all-on-all'' conjunction screening problem requires \emph{linear}
scaling: each RSO must be processed once to compute its radial bounds, and
then pairs are eliminated by simple comparison. Both the AP-filter and
SO-filter achieve this. The GEqOE-enhanced SO-filter must preserve this
property.

\subsection{Per-RSO computation (one-time cost)}

For each RSO at epoch $t_0$:
\begin{enumerate}
\item \textbf{Osculating $\to$ mean conversion}: Evaluate the GEqOE inverse
map~\eqref{eq:inverse_map}. Cost: one evaluation of the short-period map
($\sim\!100$ floating-point operations for the $J_2$-only path,
$\sim\!500$ for $J_2$--$J_5$).
\item \textbf{Mean propagation}: Integrate the mean ODE over the screening
window (5 days). For pure $J_2$: analytic (uniform rotations). For
$J_2$--$J_5$: one cheap ODE integration (the mean ODE is 5-dimensional and
slowly varying; typical step sizes are $\sim\!0.1$~days).
\item \textbf{Radial bounds}: At each evaluation epoch within the screening
window, compute $r_{\max}$ and $r_{\min}$ via Strategy~A
(Section~\ref{sec:radius_replacement}): sweep $\bar K \in [0, 2\pi]$ with
$N_{\rm grid}$ evaluations of the forward map.
\end{enumerate}

Total per-RSO cost: comparable to the current SO-filter. The dominant cost
is the $N_{\rm grid}$ forward-map evaluations, which replace the quartic solve
plus Hessian evaluation.

\subsection{Pair comparison (unchanged)}

The pair comparison step is identical to the current SO-filter: compare
$r_{\max,i} + b_i$ against $r_{\min,j} - b_j$ for all pairs $(i,j)$.
The buffers $b_i$ are smaller (Section~\ref{sec:buffer}), so \emph{more pairs
are eliminated}, reducing the downstream workload.

\subsection{Implementation options}

\paragraph{Pure Python/MATLAB prototype.}
The algebraic short-period map can be implemented directly in Python or
MATLAB. This is the fastest path to validation and comparison with the
existing SO-filter implementation (which runs in MATLAB R2021b).

\paragraph{Compiled evaluation via heyoka cfunc.}
For production use, the GEqOE short-period map expressions can be compiled
into optimized evaluation kernels using heyoka's \texttt{cfunc} facility (see
the companion note \texttt{geqoe\_mean\_taylor\_prop}). This provides:
\begin{itemize}
\item SIMD vectorization for batch evaluation across the $K$-grid;
\item automatic differentiation for gradient-based extremum finding
(avoiding the grid sweep entirely);
\item compilation once, evaluation for all RSOs with different mean states
passed as parameters.
\end{itemize}

\paragraph{Taylor polynomial maps.}
For ensemble screening (multiple epochs, parameter sensitivity), the mean ODE
can be integrated using heyoka's Taylor polynomial machinery to produce
cached polynomial maps $\bar{\bm y}(t; \bar{\bm y}_0)$ valid over the entire
5-day window. See the companion note \texttt{geqoe\_mean\_taylor\_prop},
Path~C.

%% ====================================================================
\section{Validation Plan}
\label{sec:validation}

A rigorous validation should follow the same methodology as
Rivero et al.~\cite{rivero2025}:

\subsection{Test case validation}

Reproduce the three test orbits of Table~1 in~\cite{rivero2025} (NORAD
47961, 41460, 43518) and compare:
\begin{itemize}
\item GEqOE-SO radial bounds vs.\ high-fidelity propagated bounds (analogous
to Fig.~4 of~\cite{rivero2025});
\item GEqOE-SO model error $\varepsilon$ vs.\ SO-filter model error
(Fig.~6);
\item radial bound evolution over one apsidal precession period.
\end{itemize}

\subsection{Population-level validation}

Using the same 16,972-orbit dataset:
\begin{enumerate}
\item Compute GEqOE-SO radial bounds for all orbits over a 5-day window.
\item Compare with high-fidelity numerically propagated bounds.
\item Compute model error distribution (analogous to Fig.~6).
\item Determine category-specific correction buffers (analogous to Table~3).
\item Run the full pair comparison ($\sim\!144$~million pairs) and compute
$\rho_{\rm FP}$, $\rho_{\rm FN}$, and $\eta$ (analogous to Tables~2 and~4).
\end{enumerate}

\subsection{Extended eccentricity validation}

Extend the dataset to include orbits with $0.1 < e < 0.8$:
\begin{enumerate}
\item Select high-eccentricity orbits from the full TLE catalogue.
\item Propagate with the high-fidelity model.
\item Compute GEqOE-SO bounds and model errors.
\item Determine buffers for the new eccentricity categories.
\end{enumerate}

This constitutes the primary novel contribution of the GEqOE enhancement:
analytical first-stage filtering for the currently excluded population.

%% ====================================================================
\section{Summary of Improvements}
\label{sec:summary}

Table~\ref{tab:summary} provides an overview of the component-by-component
improvements.

\begin{table}[h]
\centering
\caption{Summary of GEqOE enhancements to the SO-filter.}
\label{tab:summary}
\small
\begin{tabular}{p{3.2cm}p{4.5cm}p{4.5cm}}
\toprule
\textbf{Component} & \textbf{Current SO-filter} & \textbf{GEqOE-SO} \\
\midrule
Osc-to-mean &
Kozai--Lyddane, $J_2$-only, singular at $i_c$ &
GEqOE inverse map, $J_2$--$J_5$, singularity-free \\[4pt]
Mean evolution &
Cook's circle, 1st-order-in-$e$, $J_2$ rate &
GEqOE mean ODE, exact-in-$e$, $J_2$--$J_5$ \\[4pt]
Radius model &
1st-order-in-$e$, $J_2$-only polar equation &
Exact-in-$e$ osculating reconstruction \\[4pt]
Eccentricity range &
$e < 0.1$ &
$e < 0.8+$ (full catalogue) \\[4pt]
Correction buffer &
0.5--2.5~km &
0.2--0.7~km (est., low-$e$) \\[4pt]
Drag model &
Billik altitude decay, circular only &
King-Hele secular rates, arbitrary $e$ \\[4pt]
Computational cost &
$\sim\!8$~min for 16,972 orbits &
Comparable ($\sim\!10$~min est.) \\
\bottomrule
\end{tabular}
\end{table}

%% ====================================================================
\section{Collaboration Path}
\label{sec:collaboration}

The integration of the GEqOE theory into the SO-filter framework suggests a
natural collaboration between the analytical orbit theory development (this
work and the master note) and the conjunction analysis group (Rivero,
Bombardelli, Vazquez). A concrete path forward:

\begin{enumerate}
\item \textbf{Phase~1: Prototype validation.} Implement the GEqOE inverse map
and forward map in MATLAB (to interface with the existing SO-filter code).
Validate on the three test orbits of~\cite{rivero2025}. Deliverable: error
comparison plots analogous to Fig.~4.

\item \textbf{Phase~2: Population benchmark.} Run the GEqOE-SO filter on the
full 16,972-orbit dataset. Compute buffers and filter metrics. Deliverable:
updated Tables~2--4 with GEqOE-SO column.

\item \textbf{Phase~3: Eccentricity extension.} Extend the dataset to
$e > 0.1$. Compute GEqOE-SO bounds and buffers for the high-$e$ population.
Deliverable: first analytical filter results for GTO/Molniya-class orbits.

\item \textbf{Phase~4: Drag integration.} Replace the Billik formula with
the GEqOE drag extension for low-altitude orbits. Deliverable: updated
Table~6 with GEqOE-SO column.

\item \textbf{Phase~5: Production implementation.} Port to compiled heyoka
cfunc kernels for production-speed evaluation. Deliverable: timing benchmarks
and integration with space traffic management tools.
\end{enumerate}

Each phase produces a self-contained, publishable result:
\begin{itemize}
\item Phase~1--2: journal paper comparing GEqOE-SO vs.\ SO-filter on the
standard benchmark;
\item Phase~3: conference paper on high-eccentricity conjunction filtering;
\item Phase~4--5: applications paper on integrated drag + zonal conjunction
screening.
\end{itemize}

\begin{thebibliography}{99}
\bibitem{rivero2025}
Rivero, A.S., Bombardelli, C., Vazquez, R.,
``Short-term space occupancy and conjunction filter,''
\emph{Advances in Space Research}, 76, 2354--2372, 2025.

\bibitem{cook1966}
Cook, G.,
``Perturbations of near-circular orbits by the Earth's gravitational
potential,''
\emph{Planetary and Space Science}, 14(5), 433--444, 1966.

\bibitem{kozai1959}
Kozai, Y.,
``The motion of a close earth satellite,''
\emph{The Astronomical Journal}, 64, 367--377, 1959.

\bibitem{lyddane1963}
Lyddane, R.H.,
``Small eccentricities or inclinations in the Brouwer theory of the
artificial satellite,''
\emph{The Astronomical Journal}, 68, 555--558, 1963.

\bibitem{bau2021}
Ba\`u, G., Bombardelli, C., Peláez, J.,
``A generalization of the equinoctial orbital elements,''
\emph{Celestial Mechanics and Dynamical Astronomy}, 133, Article 50, 2021.

\bibitem{bombardelli2021}
Bombardelli, C., Falco, G., Amato, D., et al.,
``Space occupancy in low-earth orbit,''
\emph{Journal of Guidance, Control, and Dynamics}, 44(4), 684--700, 2021.
\end{thebibliography}

\end{document}
